\documentclass[a4paper,11pt]{article} % Die Angaben in eckigen Klammern sind optional (Standardschriftgröße ist 10pt)
%%%%%%%%%%%%%%%%%%%%%%%%%%%%%%%%%%%%%%%%%%%%%%%%%%%%%%%%%%%%%%%%%%%%%%%%%%%%%%%%%%%%%%%%%%%%%%%%

%
%%%%%%%%%%%%%%%%%%%%%%%%%%%%%%%%%%%%%%%%%%%%%%%%%%%%%%%%%%%%%%%%%%%%%%%%%%%%%%%%%%%%%%%%%%%%%%%%
\usepackage[english,ngerman]{babel} % So ist deutsch die Standardsprache, englisch kann eingeschaltet werden.
\usepackage[utf8]{inputenc} % Für Umlaute
\usepackage[T1]{fontenc} % Für bessere Silbentrennung
\usepackage{lmodern} % Für eine weniger pixelige Schriftart
\usepackage{graphicx} % Für Bilder (png, pdf, jpg,...)

%-------------Mathezeugs--------------
\usepackage{amssymb}
\usepackage{amsthm} % Für mehr Funktionen bei den theorem-Umgebungen
\usepackage{mathtools} % Eine Weiterentwicklung von amsmath
\usepackage{hyperref}
%
%-------Abstände etc.--------------
\usepackage{geometry} % Zur individuellen Anpassung der Seitenränder, z.B.:
\geometry{top=15mm, bottom=15mm, left=20mm, right=15mm}
%
\setlength{\parindent}{0ex} % Kein Einrücken bei einem Absatz
\renewcommand{\arraystretch}{1.5}   % Zeilenbreite bei Tabellen
\pagestyle{empty}                   % keine Seitennummern
%
%----------Kurzbefehle für Mathe-Buchstaben -------------------
\newcommand{\NN}{\mathbb N}
\newcommand{\ZZ}{\mathbb Z}
\newcommand{\QQ}{\mathbb Q}
\newcommand{\RR}{\mathbb R}
\DeclareMathOperator{\ggt}{ggT}
%
%------------Theorem-Definitionen----------
% braucht man für Seminararbeiten o.Ä„.
\newtheoremstyle{costum}
  {\topsep}   % ABOVESPACE
  {\topsep}   % BELOWSPACE
  {\normalfont}  % BODYFONT
  {0pt}       % INDENT (empty value is the same as 0pt)
  {\itshape} % HEADFONT
  {.}         % HEADPUNCT
  {5pt plus 1pt minus 1pt} % HEADSPACE
  {\thmname{#1}\thmnumber{ #2}\thmnote{\textnormal{ \textbf{#3}}}}          % CUSTOM-HEAD-SPEC
\theoremstyle{costum} % Die folgenden Umgebungen werden damit nicht kursiv
\newtheorem*{satz}{Satz}%[section] % Die Nummerierung für Sätze richtet sich so nach Kapiteln
% Das [satz] in den folgenden Definitionen bedeutet, dass alle diese Theoreme fortlaufend mit den Sätzen nummeriert werden (Also "Satz 1.1, Bemerkung 1.2, Satz 1.3" statt "Satz 1.1, Bemerkung 1.1, Satz 1.2" etc.)
% \newtheorem{prop}[satz]{Proposition}
% \newtheorem{lem}[satz]{Lemma}
\newtheorem*{kor}{Korollar}
\newtheorem*{bem}{Bem}
\newtheorem*{defi}{Def}
\newtheorem*{bsp}{Bsp}
\newtheorem*{lemma}{Lemma}

\newcommand{\gdw}{\Leftrightarrow}

%%%%%%%%%%%%%%%%%%%%%%%%%%%%%%%%%%%%%%%%%%%%%%%%%%%%%%%%%%%%%%%%%%%%%%%%%%%%%%%%%%%%%%%%%%%%%%%%
\usepackage{enumitem}
\setlist{nolistsep}
%
\begin{document}

\raisebox{-5mm}{\includegraphics[width=30mm]{Bilder/kit_logo_de_color.jpg}}
%
\hfill \parbox{24mm}
{
Malte Vo\ss\\
Larissa Bihr\\
}

\rule{\textwidth}{1pt}                                   % Linie über die ganze Textbreite
%
\begin{center}
\textbf{
Vokabeln der Algebra und Zahlentheorie \\[1ex] % Zeilenumbruch mit Abstand
}
{
\today, \url{https://github.com/vossmalte}
}
\end{center}
%
\rule{\textwidth}{1pt}\\                                 % Linie über die ganze Textbreite
%%%%%%%%%%%%%%%%%%%%%%%%%%%%%%%%%%%%%%%%%%%%%%%%%%%%%%%%%%%%%%%%%%%%%%%%%%%%%%%%%%%%%%%%%%%%%%%%
\section{Teilbarkeit}

\defi{} $d$ teilt $n$ $\Leftrightarrow$ $\exists t: t \cdot d = n$ \\
Begriffe: Teiler, Vielfaches, ggT, kgV

\bem{ggT} $ggT(0,m)=m$, $ggT(m,n) = ggT(m,n-m)$

\defi{Euklidischer Algorithmus.} $a\stackrel{!}{<}b$. Solange $b-a>0$: $b:= b-a$

\bem{} mit dem eukl. Alg. kann man auch den ggT als Linearkombination darstellen.

\defi{Kongruenz} $k\equiv l (\mod n)$ $\Leftrightarrow n | k-l$

\bem{}
\begin{itemize}
    \item $g=ggT(n,m)$ $ \Rightarrow ggT(\frac{n}{g},\frac{m}{g}) = 1$
    \item $m,n$ teilerfremd und $m \mid nu$ $\Leftrightarrow m\mid u$
    \item $kgV(m,n) \cdot ggT(m,n) = m \cdot n$
\end{itemize}

\defi{Primzahlen} ...

\bem{Fundamentalsatz der Arithmetik:} Jede Zahl hat eine Primfaktorzerlegung

\defi{$p$-adische Bewertung.} $v_p(k)$ ist die höchste Potenz von $p$, die $k$ teilt: $p^{v_p(k)} \mid k$

\bem{}
\begin{itemize}
    \item $b \mid a \gdw \forall p \in \mathbb{P}: v_p(b)\leq v_p(a)$
    \item ggT von $a$ und $b$ ist $\prod_{p\in\mathbb{P}}p^{\min(v_p(a),v_p(b)}$
    \item kgV von $a$ und $b$ ist $\prod_{p\in\mathbb{P}}p^{\max(v_p(a),v_p(b)}$
\end{itemize}

\satz{kl. von Fermat:} $p\in\mathbb{P}$ und $c\in\mathbb{Z}$ $\Rightarrow p \mid c^p-c$ { \footnotesize (hiermit kann prim wiederlegt werden)}

\bem{Lücken:} $k\in\mathbb{N}$. Zwischen $M:=(k+2)!+2$ und $M+k$ liegt keine Primzahl


\section{Gruppen}

\subsection{Magmen}

\defi{Magma} $(M,*)$ ist eine Menge mit einer Verknüpfung $*:M\times M \to M$
\begin{itemize}
    \item ist die Verknüpfung assoziativ, so nennt man das Magma Halbgruppe
    \item Eine Halbgruppe mit neutralem Element heißt Monoid
    \item es gibt auch kommutative Magmen
    \item Bsp: Abb(D,D) mit Komposition ist ein Magma. Leeres Magma. Triviales Magma.
    \item $U\subset M$ heißt Untermagma wenn $U*U\subset U$. Untermagmen kann man schneiden und diese bleiben Untermagmen
    \item Sei $X\subset M$, das Magmenerzeugnis $\langle X\rangle$ von $X$ ist der Schnitt alles Untermagmen, die $X$ enthalten.
    \item für assoz. Magmen gilt: $\langle X\rangle = \cup_{n\in\mathbb{N}}X_n$ mit $X_1=X, X_{n+1}=X*X_n$ 
\end{itemize}

\subsection{Der Gruppenbegriff}

\defi{Gruppe} 
\begin{itemize}
    \item Sei $(M,*)$ ein Monoid mit neutralem Element $e$. Ein Element $x$ heißt invertierbar, wenn ein $y \epsilon M$ exisitiert, sodass $x*y=y*x=e$. Bezeichne mit $M^x$ die Menge aller invertierbaren Elemente in M.
    \item Eine Gruppe ist ein Monoid, in dem jedes Element invertierbar ist.
    \item Sei $(M,*)$ eine Gruppe. $(M,*)$ ist kommutativ/abelsch wenn sie als Magma kommutativ ist.
\end{itemize}

\defi{Untergruppe}
Sei $(G,*)$ eine Gruppe. Dann ist eine Untergruppe von G ein nichtleeres Untermagma U, das unter der Inversenbildung abgeschlossen ist. \\
Schreibe: $U \leq \ G$ 
\begin{itemize}
    \item Die Gruppe $(\mathbb{Z},+)$ ist eine Untergruppe von $(\mathbb{Q},+)$.
    \item Die Gruppe der ganzen Zahlen $\mathbb{Z}$ mit der Addition als Verknüpfung. Die triviale Untergruppe ist die ${0}$. AUßerdem gilt für jede natürliche Zahl n die Teilmenge \\
    $n \mathbb{Z}:={nk \vert k \epsilon \mathbb{Z}}$ \\
    Wir halten fest: Die Untergruppen von $\mathbb{Z}$ sind genau die Mengen $n \mathbb{Z}$ mit $n \epsilon \mathbb{N}_{0}$
\end{itemize}

\bem{} Durchschnitt von Untergruppen: Sei G Gruppe, I nichtleere Menge und für jedes $i \epsilon I$ gibt es eine Untergruppe $U_i$ von G geben. Dann ist $\cap_{i \epsilon I} U_i$ eine Untergruppe von G.

\defi{Gruppenerzeugnis}: Sei M Teilmenge der Gruppe G, sei I die Menge aller Untergruppen von G, die M enthalten. \\
$\langle M \rangle:=\cap_{i \epsilon I} i$ 
ist Gruppenerzeugnis von M oder die von M erzeugte Untergruppe von G.
\begin{itemize}
    \item Es ist die kleinste Untergruppe von G, die M enthält. 
\end{itemize}

\defi{zyklisch}: Eine Gruppe G heißt zyklisch, wenn es ein Element $a\epsilon G$, sodass $G=\langle{a}\rangle$

\defi{Ordnung} Die Mächtigkeit/Kardinalität einer Gruppe nennt man auch Ordnung. Die Ordnung eines Elemtens $g\epsilon G$ ist definiert als die Ordnung der von g erzeugten Untergruppe.
\begin{itemize}
    \item Wenn $g\epsilon G$ endliche Ordnung, dann ist diese gleich der kleinsten natürlichen Zahl k, für die $g_k=e_G$ gilt.
\end{itemize}

\bem{} Satz von Lagrange: Sei G endl. Gruppe und H Untergruppe von G. Dann gilt: $\#H\mid\#G$

\defi{Index}: Wenn $H \leq \ G$ zwei Gruppen, dann heißt die Anzahl der Äquivalenzklassen der Index von H in G. Schreibe $(G:H)$. \\
Es gilt demnach für alle Gruppen $(G:H)=\frac{\#G}{\#H}$.

\subsection{Homomorphismen zwischen Gruppen}
\defi{Gruppenhomomorphismus}: Sei $(G,*)$ und $(H,\cdot)$ zwei Gruppen. Ein Homomorphismus von G nach H ist eine Abbildung $f:G\rightarrow H$, für die gilt:

\section{Teilbarkeit und Primelemente}
\subsection{}
\defi{Assoziert} sind $a$ und $b$ wenn $\exists e \in R^\times: b=a\cdot e$

\bem{Ordnung} entsteht durch Teilbarkeit: $aR^\times \leq bR^\times \gdw a|b$
\bem{teilerfremd} sind $a$ und $b$, wenn die einzigen Teiler Einheiten sind.

\defi{Hauptideal} $I$ ist ein Ideal, für das gilt: $I=Rg=(g)$

\defi{Hauptidealring} ist ein nullteilerfreier kommutativer Ring, in dem jedes Ideal Hauptideal ist.

\bem{ggT:} für $I=Rg={ax+by \mid x,y \in R}$ ist $g=\ggt(a,b)$

\kor{Erzeuger} $Rg=Rh \gdw g,h$ sind assoziert. 

\defi{Euklidischer Ring} ist ein nullteilerfreier kommutativer Ring $R$ mit $\gamma:R\to \NN_0$, wobei $\gamma(r)=0 \gdw r=0$ und $\forall a,b \neq 0 \exists c: \gamma(a-bc)<\gamma(b)$

% weiter bei 5.1.10
\bem{Euklidische Ringe} sind Hauptidealringe

\bem{Chinesischer Restsatz:} $R/(Rrs) \cong R/(Rr) \times R/(Rs)$

\subsection{Arithmetik in Hauptidealringen}

\defi{$m$ irreduzibel:} $\gdw \forall a,b: m=ab \Rightarrow a \text{ oder } b \in R^\times$
\defi{$p$ prim:} $\gdw \forall a,b: p \mid ab \Rightarrow p|a \text{ oder } p \mid b$

\kor{Primelemente} von 0 verschieden sind immer irreduzibel
\kor{Im Hauptidealring} sind irreduzible Elemente prim.

\satz{Primzerlegung} $\mathbb{P}_R$ ist Vertretersystem der Primelemente. Dann ist $\forall r\in R$ ein eindeutiges Produkt assoziert aus Elementen in $\mathbb{P}_R$.

\bem{Restklassenkörper} $R/Rg$ ist Körper $\gdw g$ irreduzibel

\defi{Primideal:} $xy\in I \Rightarrow x\in I \text{ oder } y \in I$






\newpage
\section{Übersicht über Strukturen und Homomorphismen}
\begin{tabular}{p{.4\textwidth}|p{.4\textwidth}}
    Struktur &  Homomorphismus\\
    \hline\hline
    Magma: $(M,*)$
     & $\phi(m*n)=\phi(m)*\phi(n)$ \\
     \hline
    Halbgruppe ist assoziatives Magma. & \\
    \hline
    Monoid ist Halbgruppe mit neutralem Element. & 
     $\phi(e)=e$
     
    \\\hline
    Gruppe ist Monoid und alle Elemente sind invertierbar. &
     $\phi(x^{-1})=\phi(x)^{-1}$ \newline
     $\text{Kern}\phi=\phi^{-1}(\{e\})$ ist Untergruppe \newline
     $\phi(G)$ ist Untergruppe \newline
     $\phi$ injektiv $\gdw \phi^{-1}(\{e\})=\{e\}$ \newline
    
    \\\hline
    Normalteiler $N$ von $G$: $\forall g: gNg^{-1}=N$ &
     Jeder Kern ist Normalteiler
    
    \\\hline
    Faktorgruppe $G/N$ mit $(gN)\cdot(hN):=ghN$ &
     $\text{Kern}\phi=N \Rightarrow gN \mapsto \phi(g)$ ist Isomorphismus \newline
     $H$ Untergruppe: $H/(N\cap H) \cong (HN)/N$
    
    \\\hline
    Gruppenoperationen $G\times M \to M$ mit $e\cdot m = m$ und $(gh)m = g(hm)$ \newline Bahn, Stabilisator, Fixpunkt, transitiv &
    
    \\\hline
    Ring $(R,+,\cdot)$ ist bzgl $+$ abelsche Gruppe und bzgl $\cdot$ Monoid mit Distributivgesetzen.
    &
     $\text{Kern}\phi=\phi^{-1}({0})$ \newline
     $\phi(R^\times)\subset S^\times$
    \\\hline
    Integritätsbereich ist kommutativer, nullteilerfreier Ring.
    \\\hline
    Körper ist kommutativer Ring und $R^\times = R\setminus {0}$
    \\\hline
    Ideal: $\forall x\in I, r\in R: xr,rx \in I$ ist ein Normalteiler bzgl $+$
    \\\hline 
    Faktorring $R/I$ &
     $I\subset \text{Kern}\phi \Rightarrow \exists \Tilde{\phi}: \phi=\Tilde{\phi}\circ\phi$
    \\\hline
    $R$-Modul $M$ ist abelsche Gruppe mit \newline
    $(r+s)m=rm+sm$ \newline
    $r(m+n)=rm+rn$ \newline
    $(rs)m=r(sm)$ \newline
    $1m=m$ &
     Module über Körpern sind Verktorräume
    \\\hline
    $R$-Algebra A mit $\sigma:R\to A$ mit $\sigma(r)a=a\sigma(r)$
    
\end{tabular}
\end{document}
